\documentclass[11pt,a4paper]{article}

\usepackage[dutch]{babel}
\usepackage[utf8]{inputenc}
\usepackage[T1]{fontenc}
\usepackage{lmodern}
\usepackage[margin=2.5cm]{geometry}
\usepackage{graphicx}
\usepackage{booktabs}
\usepackage{amsmath}
\usepackage{amssymb}
\usepackage{caption}
\usepackage{subcaption}
\usepackage{float}
\usepackage{multirow}
\usepackage{hyperref}
\usepackage{xcolor}
\usepackage{listings}
\usepackage{algorithm}
\usepackage{algpseudocode}

\hypersetup{
    colorlinks=true,
    linkcolor=blue!60!black,
    citecolor=blue!60!black,
    urlcolor=blue!60!black
}

\title{\textbf{Tabular Reinforcement Learning op Taxi-v3} \\
       \large Een vergelijking van Q-learning en SARSA tegen baselines}
\author{Karol Gackowski \and Joep Houweling \and Louis Brouwer \and Valeriy Komarov}\date{\today}

\begin{document}
\maketitle

\begin{abstract}
\noindent
We onderzoeken in welke mate twee klassieke tabulaire Reinforcement
Learning algoritmen --- Q-learning en SARSA --- de Taxi-v3 omgeving uit
Gymnasium kunnen oplossen. Beide algoritmen zijn vanaf nul ge\"implementeerd
in NumPy, getraind op 3000 episodes en vergeleken met een uniform-random
baseline en een Manhattan-shortest-path heuristiek die geen kennis heeft
van muren. Q-learning bereikt een gemiddeld rendement van $+8{,}33$ met
$100\%$ succesvolle leveringen; SARSA $-2{,}32$ met $95\%$ succes. Beide
verslaan ruimschoots de heuristiek ($-156$, $21\%$) en de random baseline
($-748$, $6\%$). Hyperparameter-sweeps over $\alpha$, $\gamma$ en het
$\varepsilon$-decay-schema laten zien dat Q-learning robuust is voor
$\alpha$ en $\gamma$, maar gevoelig is voor te trage $\varepsilon$-decay,
terwijl SARSA bij hoge $\alpha$ catastrofaal instort vanwege de
on-policy bootstrap. We koppelen deze bevindingen terug aan de theorie en
bespreken de beperkingen van tabulaire methoden.
\end{abstract}

\section{Introductie}

Reinforcement Learning (RL) bestudeert hoe een autonome agent door
interactie met een omgeving een gedragsstrategie kan leren die de
verwachte cumulatieve beloning maximaliseert
\cite{sutton_barto_2018}. In tegenstelling tot supervised learning is er
geen vooraf gegeven dataset met correcte acties; de agent ontdekt zelf
welke acties op de lange termijn lonend zijn via trial-and-error.

Voor deze portfolio-opdracht hebben wij gekozen voor de \emph{Taxi}
omgeving, oorspronkelijk geformaliseerd door
\citet{dietterich_2000} en tegenwoordig opgenomen als
\texttt{Taxi-v3} in de Gymnasium-suite \cite{towers_2024}. De agent
bestuurt een taxi op een $5 \times 5$ raster met muren, moet een
passagier ophalen op \'e\'en van vier vaste lokaties en deze afleveren op
een eveneens vooraf bepaalde bestemming. Elke stap kost \'e\'en
beloningseenheid; een illegale pickup of dropoff kost tien; een
succesvolle aflevering levert twintig op.

\paragraph{Waarom RL geschikt is.}
Drie eigenschappen maken Taxi-v3 een natuurlijk doelwit voor RL en
ongeschikt voor supervised of regelgebaseerde alternatieven:

\begin{enumerate}
  \item \textbf{Geen gelabelde dataset.} Er bestaat geen verzameling
        ``correcte'' acties per toestand om uit te leren; de agent moet
        zelf experimenteren.
  \item \textbf{Lange-termijn credit-assignment.} De $+20$ beloning bij
        aflevering moet teruggekoppeld worden naar de tientallen acties
        die er aan vooraf gingen. Een hand-gecodeerde Manhattan-heuristiek
        kan dit niet, want zij ``ziet'' geen muren en raakt vaak vast
        (zie sectie~\ref{sec:results}).
  \item \textbf{Gestructureerde sub-taken.} De optimale strategie hangt
        af van of de passagier al in de taxi zit. Een statische heuristiek
        moet deze logica expliciet coderen; RL leert dit impliciet via
        de toestandsafhankelijkheid van $Q(s,a)$.
\end{enumerate}

\paragraph{Onderzoeksvragen.} Dit rapport beantwoordt:
(i) hoe presteren Q-learning en SARSA, vanuit nul ge\"implementeerd, op
Taxi-v3 ten opzichte van een random- en een heuristische baseline?
(ii) hoe gevoelig zijn beide algoritmen voor de keuze van $\alpha$,
$\gamma$ en het $\varepsilon$-decay-schema?
(iii) welke kwalitatieve verschillen tussen on-policy (SARSA) en
off-policy (Q-learning) leren observeren we in dit domein?

\section{Probleemdefinitie en MDP}
\label{sec:mdp}

Taxi-v3 wordt gemodelleerd als een Markov Decision Process (MDP)
$\langle \mathcal{S}, \mathcal{A}, P, R, \gamma \rangle$ met de
volgende elementen:

\paragraph{Toestandsruimte $\mathcal{S}$.}
$|\mathcal{S}| = 500$ discrete toestanden, gecodeerd als
\begin{equation}
  s = \bigl((r \cdot 5 + c) \cdot 5 + p\bigr) \cdot 4 + d,
\end{equation}
waarbij $r, c \in \{0,\dots,4\}$ de positie van de taxi zijn,
$p \in \{0,\dots,4\}$ de locatie van de passagier (0=R, 1=G, 2=Y,
3=B, 4=in\_taxi) en $d \in \{0,\dots,3\}$ de bestemming.

\paragraph{Actieruimte $\mathcal{A}$.}
$|\mathcal{A}| = 6$: \{south, north, east, west, pickup, dropoff\}.
Bewegingsacties die tegen een muur of de rand bewegen, laten de agent
op zijn plaats. Pickup en dropoff buiten een geldige locatie zijn
toegestaan maar kosten beloning.

\paragraph{Belonings\-functie $R$.}
$-1$ per stap, $-10$ voor een illegale pickup of dropoff,
$+20$ voor een succesvolle aflevering (terminale beloning).

\paragraph{Discount $\gamma$.}
We gebruiken $\gamma = 0{,}99$ als basis. Omdat de episodes terminaal
zijn na succesvolle aflevering, is $\gamma$ vooral een knop voor de
``urgentie'' waarmee de agent wil dat de aflevering snel gebeurt.

\section{Methode}
\label{sec:method}

\subsection{Algoritmen}

\paragraph{Q-learning (off-policy)} \cite{watkins_dayan_1992}
schat de actie-waardefunctie $Q(s,a)$ met de update-regel
\begin{equation}
  Q(s_t, a_t) \leftarrow Q(s_t, a_t) + \alpha
    \bigl[ r_{t+1} + \gamma \max_{a'} Q(s_{t+1}, a') - Q(s_t, a_t) \bigr].
  \label{eq:qlearning}
\end{equation}
Het bootstrap-target gebruikt de \emph{greedy} actie in $s_{t+1}$,
ongeacht welke actie de gedragspolitiek (hier $\varepsilon$-greedy)
werkelijk gaat selecteren. Daarmee leert Q-learning de optimale
greedy-policy, ook als de gedragspolitiek nog veel exploratie bevat.

\paragraph{SARSA (on-policy)} \cite{rummery_niranjan_1994}
bootstrappt met de actie die de gedragspolitiek \emph{daadwerkelijk}
selecteert:
\begin{equation}
  Q(s_t, a_t) \leftarrow Q(s_t, a_t) + \alpha
    \bigl[ r_{t+1} + \gamma\, Q(s_{t+1}, a_{t+1}) - Q(s_t, a_t) \bigr].
  \label{eq:sarsa}
\end{equation}
Daardoor incorporeert SARSA de exploratie-ruis in zijn waardeschatting:
de geleerde policy is optimaal \emph{gegeven} dat de gedragspolitiek
$\varepsilon$-greedy blijft.

\paragraph{Implementatie.} Beide algoritmen zijn als puur Python
ge\"implementeerd, met de Q-tabel als een
$|\mathcal{S}| \times |\mathcal{A}|$ NumPy-array. We gebruiken
$\varepsilon$-greedy actie-selectie met willekeurige tie-breaking in de
greedy stap, en lineaire $\varepsilon$-decay van $1{,}0$ naar $0{,}05$
over 2000 episodes. Geen enkele RL-bibliotheek (Stable-Baselines, RLlib,
e.d.) is gebruikt; we voldoen daarmee aan de eis dat de algoritmen
vanaf nul ge\"implementeerd zijn.

\subsection{Baselines}

\paragraph{Random.} Uniforme keuze uit de zes acties. Dient als
ondergrens.

\paragraph{Heuristiek.} Een Manhattan-shortest-path strategie die
geen kennis heeft van muren:
\begin{itemize}
  \item Als de passagier nog niet in de taxi zit: ga naar de pickup-tegel
        via de actie die de Manhattan-afstand het meest verkleint; voer
        \texttt{pickup} uit zodra je er bent.
  \item Anders: ga naar de bestemmingstegel via dezelfde regel; voer
        \texttt{dropoff} uit zodra je er bent.
\end{itemize}
Door de muren niet te kennen, loopt de agent regelmatig vast.

\subsection{Hyperparameter-experimenten}

We voeren vier sweeps uit, telkens met vijf seeds per waarde en
3000 episodes per run:

\begin{enumerate}
  \item \textbf{Q-learning, $\alpha \in \{0{,}1, 0{,}3, 0{,}5, 0{,}9\}$}
        --- gevoeligheid voor de leersnelheid.
  \item \textbf{Q-learning, $\gamma \in \{0{,}8, 0{,}9, 0{,}95, 0{,}99\}$}
        --- gevoeligheid voor de horizon.
  \item \textbf{Q-learning, $\varepsilon$-decay-episodes
        $\in \{200, 1000, 2000, 5000\}$}
        --- gevoeligheid voor het exploratie-schema.
  \item \textbf{SARSA, $\alpha \in \{0{,}1, 0{,}3, 0{,}5, 0{,}9\}$}
        --- spiegel van sweep 1, om on-policy gedrag te isoleren.
\end{enumerate}

\subsection{Evaluatieprotocol}

Na training worden agenten ge\"evalueerd over $100$ episodes met
$\varepsilon = 0$ (puur greedy) en met seeds die niet overlappen met de
trainingsseeds. We rapporteren gemiddeld rendement, gemiddelde
episodelengte, gemiddeld aantal illegale acties en het
afleveringspercentage.

\section{Resultaten}
\label{sec:results}

\subsection{Hoofdvergelijking}

Tabel~\ref{tab:main} toont de greedy-evaluatie van alle vier de agenten
na 3000 trainings\-episodes. Beide RL-algoritmen verslaan de baselines
ruimschoots; Q-learning bereikt de optimale prestatie ($100\%$ succes,
gemiddelde episode-lengte $\approx 13$, in lijn met het theoretische
optimum van ongeveer 12--14 stappen per episode), terwijl SARSA er net
onder zit met $95\%$ succes. De wandblinde heuristiek presteert
nauwelijks beter dan random ondanks haar doelgerichte ontwerp: zonder
muur\-bewustzijn raakt zij in 79\% van de episodes verstrikt en haalt
zij de tijdslimiet.

\begin{table}[H]
\centering
\caption{Greedy evaluatie ($\varepsilon=0$) over 100 episodes.
``Mean return'' is het gemiddelde gediscounteerde rendement;
``Length'' is de gemiddelde episodelengte;
``Success'' is het percentage episodes dat eindigt in een succesvolle
aflevering.}
\label{tab:main}
\begin{tabular}{lrrrr}
\toprule
Agent & Mean return & Std & Mean length & Success rate \\
\midrule
Random            & $-748{,}47$ & $130{,}67$ & $193{,}9$ & $\phantom{0}6\%$ \\
Heuristiek        & $-155{,}74$ & $\phantom{0}85{,}86$ & $160{,}2$ & $21\%$ \\
SARSA             & $\phantom{-00}-2{,}32$ & $\phantom{0}45{,}45$ & $\phantom{0}22{,}3$ & $95\%$ \\
\textbf{Q-learning} & $\phantom{-00}+\mathbf{8{,}33}$ & $\phantom{00}\mathbf{2{,}82}$ & $\phantom{0}\mathbf{12{,}7}$ & $\mathbf{100\%}$ \\
\bottomrule
\end{tabular}
\end{table}

\subsection{Trainingscurves}

Figuur~\ref{fig:reward_curves} laat de leercurves van Q-learning en
SARSA zien. Beide algoritmen vertonen het kenmerkende ``hockey-stick''
patroon: een lange initi\"ele plateau waarin random exploratie weinig
positief signaal oplevert, gevolgd door een snelle stijging zodra de
agent de eerste succesvolle leveringen ervaart en het $+20$ signaal
zich door de Q-tabel verspreidt. Q-learning convergeert iets sneller
en eindigt op een hoger plateau dan SARSA, in lijn met het verschil in
greedy-evaluatie uit Tabel~\ref{tab:main}.

\begin{figure}[H]
\centering
\includegraphics[width=0.95\linewidth]{../results/reward_curves.png}
\caption{Trainingscurves (rolling mean over 20 episodes) voor
Q-learning en SARSA op Taxi-v3. De stippellijn geeft de theoretisch
benaderde optimale return ($\approx +8$).}
\label{fig:reward_curves}
\end{figure}

\subsection{Geleerde policies}

Figuren~\ref{fig:policy_q} en~\ref{fig:policy_sarsa} visualiseren de
greedy policies in acht (passagierslocatie, bestemming) configuraties.
Voor elke configuratie tonen we $V(s) = \max_a Q(s,a)$ als heatmap
met de greedy actie als pijl per cel; \texttt{P} en \texttt{D} markeren
respectievelijk pickup- en dropoff-acties. De donkergroene gradient
loopt logisch naar de oranje (passagier) of rode (bestemming) markering.
De Q-learning policy is over alle scenario's stabiel coherent; SARSA's
policy is grotendeels gelijk maar bevat enkele cellen met een suboptimale
actie, wat het lagere succespercentage in Tabel~\ref{tab:main} verklaart.

\begin{figure}[H]
\centering
\includegraphics[width=\linewidth]{../results/qlearning_default/policy.png}
\caption{Q-learning greedy policy in acht scenario's. Bovenste rij:
passagier wacht bij \textbf{R} (oranje cirkel). Onderste rij: passagier
zit al in de taxi. Kolommen: bestemming \textbf{R}, \textbf{G},
\textbf{Y}, \textbf{B} (rood vierkant).}
\label{fig:policy_q}
\end{figure}

\begin{figure}[H]
\centering
\includegraphics[width=\linewidth]{../results/sarsa_default/policy.png}
\caption{SARSA greedy policy, zelfde acht scenario's als
Figuur~\ref{fig:policy_q}.}
\label{fig:policy_sarsa}
\end{figure}

\subsection{Hyperparameter-sweeps}
\label{sec:sweeps}

Tabel~\ref{tab:sweeps} vat de gemiddelde return over de laatste 200
training-episodes samen voor elke sweep, gemiddeld over vijf seeds.
Figuren~\ref{fig:sweep_alpha}, \ref{fig:sweep_gamma},
\ref{fig:sweep_eps} en~\ref{fig:sweep_sarsa_alpha} tonen de
bijbehorende leercurves met $\pm 1\sigma$ banden.

\begin{table}[H]
\centering
\caption{Gemiddelde return over de laatste 200 episodes per sweep
(mean $\pm$ std over 5 seeds).}
\label{tab:sweeps}
\begin{tabular}{lll}
\toprule
Sweep & Waarde & Mean return (last 200) \\
\midrule
\multirow{4}{*}{Q-learning, $\alpha$}
  & 0{,}1 & $+5{,}59 \pm 0{,}33$ \\
  & 0{,}3 & $+5{,}91 \pm 0{,}21$ \\
  & 0{,}5 & $+5{,}53 \pm 0{,}34$ \\
  & 0{,}9 & $+5{,}62 \pm 0{,}19$ \\
\midrule
\multirow{4}{*}{Q-learning, $\gamma$}
  & 0{,}8  & $+5{,}70 \pm 0{,}32$ \\
  & 0{,}9  & $+5{,}57 \pm 0{,}36$ \\
  & 0{,}95 & $+5{,}77 \pm 0{,}36$ \\
  & 0{,}99 & $+5{,}53 \pm 0{,}34$ \\
\midrule
\multirow{4}{*}{Q-learning, $\varepsilon$-decay-eps.}
  & 200  & $+5{,}30 \pm 0{,}22$ \\
  & 1000 & $+5{,}24 \pm 0{,}30$ \\
  & 2000 & $+5{,}53 \pm 0{,}34$ \\
  & 5000 & $\mathbf{-35{,}27 \pm 1{,}18}$ \\
\midrule
\multirow{4}{*}{SARSA, $\alpha$}
  & 0{,}1 & $+4{,}94 \pm 0{,}50$ \\
  & 0{,}3 & $+5{,}10 \pm 0{,}26$ \\
  & 0{,}5 & $+3{,}07 \pm 1{,}06$ \\
  & 0{,}9 & $\mathbf{-56{,}36 \pm 6{,}46}$ \\
\bottomrule
\end{tabular}
\end{table}

\paragraph{Q-learning, $\alpha$ (Fig.~\ref{fig:sweep_alpha}).}
Q-learning is opvallend ongevoelig voor de leersnelheid: alle vier de
waarden eindigen binnen $\pm 0{,}3$ van elkaar. Dit is consistent met
de theoretische convergentiegarantie van Q-learning onder Robbins--Monro
voorwaarden \cite{sutton_barto_2018}; in de praktijk geldt dat zolang
$\alpha$ niet zo extreem laag is dat updates verwaarloosbaar worden,
de eindwaardefunctie convergeert naar dezelfde fixed point.

\paragraph{Q-learning, $\gamma$ (Fig.~\ref{fig:sweep_gamma}).}
Vergelijkbaar robuust over de geteste range. Voor Taxi-v3 zijn episodes
relatief kort (typisch $<20$ stappen na convergentie), waardoor het
verschil tussen $\gamma=0{,}8$ en $\gamma=0{,}99$ in de praktijk klein
blijft.

\paragraph{Q-learning, $\varepsilon$-decay (Fig.~\ref{fig:sweep_eps}).}
Hier vinden we de eerste duidelijke negatieve uitslag: bij een
$\varepsilon$-decay-window van 5000 episodes is $\varepsilon$ na 3000
training-episodes nog ruim boven $0{,}3$, zodat een aanzienlijk deel
van de \emph{evaluatie}-stappen in de laatste 200 episodes nog
willekeurig zijn. De return zakt dan naar $-35{,}27$. Dit illustreert
een belangrijk praktisch punt: ook al levert Q-learning een correcte
greedy-policy op, de gemeten trainings-return wordt sterk be\"invloed
door de gedragspolitiek.

\paragraph{SARSA, $\alpha$ (Fig.~\ref{fig:sweep_sarsa_alpha}).}
Hier vinden we de meest dramatische bevinding: SARSA is veel
gevoeliger voor $\alpha$ dan Q-learning. Bij $\alpha=0{,}9$ stort SARSA
in tot een gemiddelde return van $-56{,}36$ met slechts $91{,}5\%$
succes. De verklaring volgt uit vergelijking~\eqref{eq:sarsa}: bij
hoge $\alpha$ slaat een enkele $\varepsilon$-greedy random actie met
slechte uitkomst direct door in $Q(s_t, a_t)$ via de bootstrap
$Q(s_{t+1}, a_{t+1})$ van diezelfde slechte actie. Q-learning is hier
immuun voor, omdat zijn bootstrap altijd $\max_{a'} Q(s_{t+1}, a')$
gebruikt --- de slechte actie wordt simpelweg buiten beschouwing
gelaten.

\begin{figure}[H]
\centering
\includegraphics[width=0.95\linewidth]{../results/sweep_qlearning_alpha/sweep_curves.png}
\caption{Q-learning, sweep over $\alpha$. Mean $\pm 1\sigma$ over 5 seeds.}
\label{fig:sweep_alpha}
\end{figure}

\begin{figure}[H]
\centering
\includegraphics[width=0.95\linewidth]{../results/sweep_qlearning_gamma/sweep_curves.png}
\caption{Q-learning, sweep over $\gamma$.}
\label{fig:sweep_gamma}
\end{figure}

\begin{figure}[H]
\centering
\includegraphics[width=0.95\linewidth]{../results/sweep_qlearning_epsilon_decay_episodes/sweep_curves.png}
\caption{Q-learning, sweep over de lengte van het $\varepsilon$-decay-schema.
Bij 5000 episodes blijft $\varepsilon$ tijdens de gehele training te hoog,
met een dramatisch lagere trainings-return als gevolg.}
\label{fig:sweep_eps}
\end{figure}

\begin{figure}[H]
\centering
\includegraphics[width=0.95\linewidth]{../results/sweep_sarsa_alpha/sweep_curves.png}
\caption{SARSA, sweep over $\alpha$. In tegenstelling tot Q-learning
(Fig.~\ref{fig:sweep_alpha}) stort SARSA bij hoge $\alpha$ in.}
\label{fig:sweep_sarsa_alpha}
\end{figure}

\section{Discussie en reflectie}
\label{sec:discussion}

\subsection{Vergelijking met de baselines}

De grote sprong tussen heuristiek ($21\%$ succes) en RL
($95$--$100\%$) onderstreept waarom RL de juiste paradigmakeuze is voor
Taxi: de heuristiek faalt niet omdat zijn optimaliseringscriterium
verkeerd is, maar omdat zij een kritisch aspect van de omgeving (muren)
niet expliciet modelleert. Een rijkere heuristiek met muur\-bewustzijn
zou ongetwijfeld beter presteren --- maar zou ook handmatig
voor elk nieuw raster opnieuw geprogrammeerd moeten worden, terwijl
Q-learning en SARSA dezelfde implementatie kunnen toepassen op iedere
discrete MDP. Dit illustreert de generaliseerbaarheid die RL
fundamenteel onderscheidt van rule-based en supervised methoden
\cite{sutton_barto_2018}.

\subsection{Q-learning versus SARSA in dit domein}

Op een omgeving als Cliff Walking
\cite[Voorbeeld 6.6]{sutton_barto_2018} ontstaat een dramatisch
verschil tussen on-policy en off-policy leren omdat de optimale policy
direct langs een gevaarlijke afgrond loopt. Op Taxi-v3 ontbreekt zo'n
``risico-randvoorwaarde''; de optimale policy is unimodaal en SARSA en
Q-learning convergeren beide naar (vrijwel) dezelfde policy. Het
overgebleven verschil --- $5\%$ minder succes voor SARSA en een
gevoeligere $\alpha$-respons --- is precies wat de theorie voorspelt:
SARSA's on-policy bootstrap incorporeert exploratie-ruis in $Q$, wat
bij hoge $\alpha$ en hoge $\varepsilon$ destabiliseert.

\subsection{Hyperparameter-bevindingen}

Drie inzichten zijn relevant voor toekomstige toepassingen:

\begin{enumerate}
  \item \textbf{Q-learning is robuust voor $\alpha$ en $\gamma$.}
        In een tabulaire setting met een eindige toestandsruimte
        levert elke redelijke combinatie convergentie naar de optimale
        Q-functie. Hyperparameter-tuning van deze twee is dus weinig
        rendabel.
  \item \textbf{$\varepsilon$-decay is wel kritisch.} Een te
        traag schema houdt de gedragspolitiek te lang random;
        evaluaties tijdens training reflecteren dan eerder de
        gedragspolitiek dan de geleerde policy.
  \item \textbf{Algoritme $\times$ hyperparameter is geen
        productruimte.} Een waarde die voor Q-learning veilig is
        ($\alpha=0{,}9$) is voor SARSA destructief. Tuning moet per
        algoritme gebeuren.
\end{enumerate}

\subsection{Beperkingen}

\paragraph{Schaal.} Tabulaire methoden vereisen \'e\'en
$Q(s,a)$-cel per staat-actie-paar. Voor Taxi-v3 ($500 \times 6 = 3000$
cellen) is dit triviaal, maar in elke realistische omgeving met
continue of hoog-dimensionale toestanden onhaalbaar. De volgende stap
is functie-approximatie via neurale netwerken (Deep Q-Networks,
\cite{mnih_2015}), waar dezelfde Bellman-update wordt toegepast op de
gradient van een neuraal netwerk in plaats van een tabel.

\paragraph{Stochasticiteit.} Taxi-v3 is een deterministische
omgeving: elke actie heeft een vast effect. Toevoeging van stochastische
overgangen (zoals de \texttt{is\_slippery}-modus van CliffWalking) zou
de Q-learning vs SARSA contrast verder kunnen versterken en is een
natuurlijke uitbreiding voor vervolgwerk.

\paragraph{Sample-effici\"entie.} Beide algoritmen vereisen rond de
2000 episodes voor convergentie. Methoden zoals Dyna-Q
\cite[\S 8.2]{sutton_barto_2018} kunnen dit verlagen door
ervaringen te ``replayen'' in een geleerd omgevingsmodel.

\section{Conclusie}

We hebben Q-learning en SARSA vanaf nul ge\"implementeerd en toegepast op
de Taxi-v3 omgeving. Q-learning bereikt $100\%$ greedy-succes met een
gemiddelde episode-lengte dichtbij het theoretisch minimum; SARSA komt
op $95\%$. Beide algoritmen verslaan zowel een uniform-random baseline
als een muur-onbewuste Manhattan-shortest-path heuristiek met
ordegrootten. Hyperparameter-sweeps met meerdere seeds tonen aan dat
Q-learning robuust is voor $\alpha$ en $\gamma$ maar gevoelig voor te
trage $\varepsilon$-decay, terwijl SARSA dramatisch instort bij hoge
$\alpha$ --- in lijn met de theoretische voorspelling dat de on-policy
bootstrap exploratie-ruis incorporeert in de waardeschatting.

De gekozen omgeving en aanpak hebben ons in staat gesteld de centrale
RL-concepten (Markov Decision Process, beloningen, $\varepsilon$-greedy
exploratie, on-policy versus off-policy bootstrap, hyperparameter-
gevoeligheid) op een werkende implementatie te demonstreren, terwijl
de tabulaire schaal van $|S| \cdot |A| = 3000$ tegelijk laat zien
waarom function-approximation methoden zoals DQN nodig zijn voor grotere
problemen.

\section*{Reproduceerbaarheid}

Alle code, configuraties en getrainde Q-tabellen zijn beschikbaar in de
projectrepository. Trainingen worden gestuurd via YAML-configuraties in
\texttt{experiments/}; resultaten worden per run weggeschreven naar
\texttt{results/<run\_name>/log.csv}. Elke commando-regelaanroep
accepteert \texttt{--seed} voor reproduceerbaarheid; de figuren in dit
rapport zijn gegenereerd door de scripts \texttt{src/plot\_results.py},
\texttt{src/plot\_policy.py} en \texttt{src/plot\_sweep.py}. De volledige
verzameling sweeps is reproduceerbaar met:

\begin{verbatim}
python -m src.sweep --agent qlearning --param alpha \
    --values 0.1 0.3 0.5 0.9 --seeds 5 --episodes 3000
python -m src.sweep --agent qlearning --param gamma \
    --values 0.8 0.9 0.95 0.99 --seeds 5 --episodes 3000
python -m src.sweep --agent qlearning --param epsilon_decay_episodes \
    --values 200 1000 2000 5000 --seeds 5 --episodes 3000
python -m src.sweep --agent sarsa --param alpha \
    --values 0.1 0.3 0.5 0.9 --seeds 5 --episodes 3000
\end{verbatim}

\section*{Gebruik van Generative AI}

Bij het schrijven van deze rapportage en de implementatie is gebruik
gemaakt van Anthropic's Claude (model claude-opus-4-7, mei 2026) voor
sparring over de structuur van het rapport, debugging van de
trainings\-pipeline en het opstellen van de LaTeX-skelet. Alle
algoritmische beslissingen, hyperparameter-keuzes en interpretaties
zijn door de auteurs geverifieerd tegen de geciteerde bronnen.

\begin{thebibliography}{9}

\bibitem{sutton_barto_2018}
Sutton, R.~S., \& Barto, A.~G. (2018).
\emph{Reinforcement Learning: An Introduction} (2nd ed.).
MIT Press.

\bibitem{watkins_dayan_1992}
Watkins, C.~J.~C.~H., \& Dayan, P. (1992).
Q-learning.
\emph{Machine Learning}, 8(3--4), 279--292.
\url{https://doi.org/10.1007/BF00992698}

\bibitem{rummery_niranjan_1994}
Rummery, G.~A., \& Niranjan, M. (1994).
\emph{On-line Q-learning using connectionist systems}
(Technical Report CUED/F-INFENG/TR 166).
University of Cambridge, Department of Engineering.

\bibitem{dietterich_2000}
Dietterich, T.~G. (2000).
Hierarchical reinforcement learning with the MAXQ value function decomposition.
\emph{Journal of Artificial Intelligence Research}, 13, 227--303.
\url{https://doi.org/10.1613/jair.639}

\bibitem{towers_2024}
Towers, M., Kwiatkowski, A., Terry, J., et al. (2024).
Gymnasium: A standard interface for reinforcement learning environments.
\emph{arXiv:2407.17032}.
\url{https://arxiv.org/abs/2407.17032}

\bibitem{mnih_2015}
Mnih, V., Kavukcuoglu, K., Silver, D., et al. (2015).
Human-level control through deep reinforcement learning.
\emph{Nature}, 518(7540), 529--533.
\url{https://doi.org/10.1038/nature14236}

\end{thebibliography}

\end{document}
